\centerline{\bf \Huge Показательная серия задач}

\begin{flushright}
        \scriptsize{Какие два слова вводят в страх Германа?\\ Проверка корректности}
\end{flushright}

\textbf{Определение.} Пусть $(a, n)=1$. \textit{Показателем} числа $a$ по модулю $n$ называется наименьшее натуральное число $d$ такое, что $a^d \equiv 1(\bmod n)$. Обозначение $d=\operatorname{ord}_n a$.

\textbf{Определение.} Функция Эйлера $\varphi(n)$ определяется как количество натуральных чисел, не превосходящих $n$, которые взаимно просты с $n$.

\textbf{Теорема Эйлера.} Пусть $(a, n)=1$. Докажите, что $a^{\varphi(n)} \equiv 1(\bmod n)$.

\textbf{Основные свойства показателей:}

0. Показатель $ord_na$ существует тогда и только тогда, когда $(a,n) = 1$

1. Пусть $ord_n a = d$. Тогда числа $a, a^2, a^3, \ldots, a^d$ дают разные остатки при делении на $n$.

2. Из предыдущего следует, что $a^x \equiv a^y \ (mod \ n) \longleftrightarrow x \equiv y \ (mod \ d)$

3. Из предыдушего следует, что $a^k - 1 \ \vdots \ n \longleftrightarrow k \ \vdots \ d$, в частности $\phi(n) \ \vdots \ ord_na$ для любого взаимно простого $a$.

\problem{0.1}
Дано нечетное простое число $p$, а также простые числа $q$ и $r$. Известно, что $q^r+1 \ \vdots \  p$. Докажите, что либо $p-1 \ \vdots \ 2 r$, либо $q^2-1 \ \vdots \ p$.

\problem{0.2}
Найдите все натуральные $n$, для которых числа $n$ и $2^n + 1$ имеют один и тот же набор простых делителей.

\problem{0.3}
Найдите все такие простые $p,q$, что $5^p + 5^q \ \vdots \ pq$

\problem{1}
Даны натуральные числа $a, n>1$. Докажите, что для каждого нечетного простого делителя $p$ числа $a^{2^n}+1$ число $p-1$ делится на $2^{n+1}$.

\problem{2}
(a) Докажите, что если $p$ --- простое, то у числа $2^p-1$ все простые делители больше $p$.

\noindent(б) Докажите, что если $p$ --- простое, то любой простой делитель числа $2^p-1$ имеет вид $2kp+1$ для некоторого натурального $k$.

\problem{3}
 Докажите, что для любого натурального $n$ число $2^n-1$ не
делится на $n$.

\newpage

\hspace{3mm}

\problem{4}
Докажите, что все простые делители чисел Ферма $2^{2^n}+1$ имеют вид
$2^{n+1}\cdot x+1$. (Именно так Эйлер нашел простой делитель у числа
$2^{2^5}+1$.)

\problem{5}
Найдите все тройки простых чисел $(p,q,r)$, такие, что $p \mid q^r+1$, $q \mid r^p+1$, $r \mid p^q+1$.

\problem{6}
Найти все пары простых чисел $(p;q)$, таких, что $pq \mid (5^p-2^p)(5^q-2^q)$.

\problem{7}
Найти все пары простых чисел $p$, $q$, таких, что $pq \mid
2^p+2^q$.

\problem{8}
Найдите все тройки взаимно простых в совокупности натуральных чисел $(a, b, c)$ таких, что для любого натурального $n$ число $\left(a^n+b^n+c^n\right)^2$ делится на $a b+b c+c a$.

\problem{9}
Даны взаимно простые числа $a$ и $b$. Докажите, что для любого нечетного делителя $d$ числа $a^{2^n}+b^{2^n}$ число $d-1$ делится на $2^{n+1}$.

\problem{10}
Докажите, что для натурального числа $a$ найдется такое натуральное число $b>a$, что $1+2^b+3^b$ делится на $1+2^a+3^a$.